\subsection{An actual odd-deck lift and the exact remaining involution condition}

\begin{statement}{D.35}{Proved}{Odd-deck boundary lifts and genuine realization of the compact CM torus}
This continues Part D(e), the CM follow-up to LM2, manuscript
\S51.16, folio 734 verso, scan 174, and \S51.15, folios 725--727,
scans 156--160. Retain the marked curve, $K_2$, $\Lambda$, the
canonical cusp/elliptic splitting and the matrix $Z$ of D.34.
Let $\Lambda_o\simeq C_3^2$ be the unique order-nine subgroup of
$\Lambda\simeq S_3\times C_3$. Its preimage in $\mathcal P_2$ has
an order-nine complement $\Delta_9$ to $K_2$.

For every $\eta\in\mathcal O_{3,2}^{\times}$, where
$\mathcal O_{3,2}=\mathbb Z_2[Z]\simeq\mathbb Z_2[\zeta_3]$,
there is an actual continuous automorphism of the free pro-two group
$K_2$ which centralizes the conjugation action of $\Delta_9$
pointwise, preserves each of the three peripheral conjugacy classes
with the same power
\[
                  \lambda_2=N_{\mathcal O_{3,2}/\mathbb Z_2}(\eta),
                                                               \tag{D.35.1}
\]
and acts on the canonical homology splitting by
$\lambda_2 I_{\rm cusp}\oplus\eta_{\rm ell}$.
Here the displayed CM identification sends $\zeta_3^2$ to $Z$,
because $t^2=\delta^2$ in D.34; all evaluations below use that map.
In particular the norm-one unit
\[
 \eta_*={-3-8\zeta_3^2\over7}
        ={1-2\zeta_3^2\over1-2\zeta_3},\qquad
 \eta_*(Z)={-3I-8Z\over7}                            \tag{D.35.2}
\]
has an actual such lift with $\lambda_2=1$ and with the compact
Tate-module reduction $M_{16}$ of D.33--D.34. This assertion retains
that Tate component; equality with the complete D.33 finite action
is not asserted.

Here is the exact residual group. Let $\mathcal O_{\partial}$ be
the subgroup of $\operatorname{Out}(K_2)$ preserving each cusp
conjugacy class with a common unit power. Set
\[
 C_o=C_{\mathcal O_{\partial}}(\Lambda_o),\qquad
 \iota(a)=\bar s a\bar s^{-1}\quad(a\in C_o).
                                                               \tag{D.35.3}
\]
Then $\iota$ is an involution and
\[
 C_{\mathcal O_{\partial}}(\Lambda)=C_o^{\iota},
 \qquad C_o\longrightarrow\mathcal O_{3,2}^{\times}
                  \text{ is onto}.                  \tag{D.35.4}
\]
The involution induces the identity on this torus target. In fact the
fixed subgroup also maps onto the entire torus:
\[
 C_o^{\iota}\longrightarrow\mathcal O_{3,2}^{\times}
                  \text{ is onto}.                  \tag{D.35.4a}
\]
This stronger assertion is realized by the actual Galois action on
$\mathcal E_0^\circ$ over $\mathbb Q(\zeta_3)$, including the
norm-one unit (D.35.2). The exact
relation to the full relative peripheral source is
\[
 1\longrightarrow\langle\operatorname{Int}(c_0)\rangle
 \longrightarrow\operatorname{Aut}_{Q_0,t}(\mathcal P_2)
 \longrightarrow C_o^{\iota}\longrightarrow1,
 \qquad\langle\operatorname{Int}(c_0)\rangle\simeq\mathbb Z_2.
                                                               \tag{D.35.5}
\]
The middle group acts identically on $Q_0$ and sends
$t$ to $t^{(\lambda_2,1)}$, with the full two-component exponent
of D.34. The remaining condition retains the prescribed finite action,
not only its compact Tate image. More precisely put
$R=T_3=\mathcal P_2/D_9(K_2)$ and let $\alpha_{33,R}$ be the action
retained by D.33. Reduction defines an actual homomorphism
\[
 C_o^{\iota}\longrightarrow
 \operatorname{Aut}_{Q_0,t}(R)/\langle\operatorname{Int}(c_0)\rangle.
                                                               \tag{D.35.5a}
\]
The fibre over $[\alpha_{33,R}]$ has precisely the outer classes of
normalized relative lifts of that recorded finite action. D.45 and its source-transport supplement prove that this
particular fibre is empty. Realization of the full missing CM pair
by another finite action remains undecided. Equations (D.35.3)--(D.35.5a) retain both the order-two
deck symmetry and the specified finite coordinates.
\end{statement}
\begin{proof}
We first give the marked surface comparison and the actual boundary
automorphisms used in the lift. In the original free basis of D.34
put
\[
 a=a_1,\quad b=a_2^{-1},\quad
 c=a_1a_4^{-1}a_3a_1^{-1},\quad
 d=a_1a_3^{-1}a_1a_2a_1^{-1}.
                                                               \tag{D.35.6}
\]
The inverse change of basis is
\[
 a_1=a,\quad a_2=b^{-1},\quad
 a_3=ab^{-1}a^{-1}d^{-1}a,\quad
 a_4=ab^{-1}a^{-1}d^{-1}c^{-1}a.                      \tag{D.35.7}
\]
Direct free reduction verifies both compositions. It also gives
\[
 c_0=cd[a,b],\qquad c_1=b^{-1}a^{-1}d^{-1}ab,
 \qquad c_2=a^{-1}c^{-1}a.                            \tag{D.35.8}
\]
Thus the first boundary is conjugate to $R=[a,b]cd$, and the other
two are inverse conjugates of $d,c$ respectively. These exact words
identify the original three peripheral classes with the standard
surface pair; inversion of the latter two generators does not change
their common power under any automorphism.

Let $\mathcal A_{\partial}$ be the closed subgroup of
$\operatorname{Aut}(K_2)$ preserving the three peripheral classes
individually with a common power, and let
$\mathcal A_{\partial}^{\rm perm}$ allow their permutation.
Their homology maps are continuous. In a basis with the cusp lattice
first and a chosen rank-two handle section last, the image of the
pure subgroup is exactly
\[
 \left\{\begin{pmatrix}\lambda I_2&B\\0&M\end{pmatrix}:
 \lambda\in\mathbb Z_2^\times,\quad B\in M_2(\mathbb Z_2),
 \quad M\in\operatorname{GL}_2(\mathbb Z_2),\quad\det M=\lambda
 \right\}.                                          \tag{D.35.9}
\]
Necessity follows from the peripheral homology and the relative
fundamental-class calculation in D.34. We prove surjectivity using
actual automorphisms, rather than lifting a matrix formally.

In the basis $(a,b,c,d)$ define
\[
 \begin{array}{c|cccc}
 &a&b&c&d\\\hline
 U&a&ba&c&d\\
 L&ab&b&c&d\\
 H&a&b&d&d^{-1}cd\\
 \theta&ca&b&bcb^{-1}&cdc^{-1}.
 \end{array}                                         \tag{D.35.10}
\]
Each row is an automorphism of the discrete free group and hence of
its free pro-two completion. For $U,L$ the inverse replaces the
extra handle letter by its inverse; for $H$ the inverse sends
$(c,d)$ to $(cdc^{-1},c)$. The inverse of $\theta$ is
\[
 a\mapsto b^{-1}c^{-1}ba,\quad b\mapsto b,\quad
 c\mapsto b^{-1}cb,\quad
 d\mapsto b^{-1}c^{-1}bdb^{-1}cb.                      \tag{D.35.11}
\]
Free reduction verifies both inverse compositions in all four cases.
The first three rows fix $R$, while $\theta(R)=cRc^{-1}$.
The row $H$ interchanges the $c,d$ peripheral classes; $U,L,\theta$
are pure. Thus all the stated homology operations come from the
specified peripheral automorphism groups.

The handle actions of $U,L$ generate a dense subgroup of
$\operatorname{SL}_2(\mathbb Z_2)$. Indeed their powers supply both
elementary matrices with every two-adic coefficient in the closure,
and elementary row reduction over $\mathbb Z_2$ generates
$\operatorname{SL}_2(\mathbb Z_2)$. Compactness of the source makes
the homology image closed. Conjugating $\theta$ by $H$ supplies the
other cusp row, and conjugating these two shears by the handle rotation
$UL^{-1}U$ supplies the other handle column. The four resulting
integral shear vectors have determinant $\pm1$ in $M_2(\mathbb Z)$.
Their powers and products therefore give every $B\in M_2(\mathbb Z_2)$.
All these shears are pure, including those obtained by conjugating
with $H$.

Finally every $\lambda\in\mathbb Z_2^\times$ occurs as the
cyclotomic character of an actual element of $G_{\mathbb Q}$.
Use its action on the geometric pro-two fundamental group of
$\mathcal E_0^\circ/\mathbb Q$, based at the rational tangent with
local parameter $w$ at $P_0$. All three removed points are rational,
so this is a pure peripheral automorphism with common power
$\lambda$. This assertion follows directly by acting on finite
\'etale covers and their tangential fibres; the action on a punctured
formal parameter is the action on its compatible $2^n$-th roots.
Compose such an actual automorphism with the determinant-one handle
operations and the four shears already constructed. This adjusts its
handle matrix and off-diagonal block to any prescribed pair in
(D.35.9), proving the claimed homology surjectivity.

We record why the needed correction kernel is pro-two. For a finitely
generated free pro-two group, the kernel of its automorphism action
on the Frattini quotient is pro-two. One proof uses the characteristic
lower exponent-two central series: after the action on its first
quotient is fixed, each successive automorphism kernel embeds into
an additive space of homomorphisms to the next elementary abelian
central layer. The finite quotients are consequently two-groups,
and the inverse limit is pro-two. The kernel of integral homology
on $\mathcal A_{\partial}^{\rm perm}$ is a closed subgroup of this
kernel. It fixes all peripheral labels, since their three distinct
nonzero mod-two homology vectors are fixed, and its common multiplier
is one.

The coprime complement facts needed here admit the same elementary
inverse-limit proof. In a finite extension with a normal two-group
and odd quotient, successive elementary abelian chief factors have
vanishing positive cohomology for the odd quotient, by averaging
with its invertible order. Induction supplies a complement and shows
that any two complements are conjugate by the normal two-group.
Finite solutions form nonempty compact sets; their inverse limit
supplies a complement and conjugacy for a normal pro-two kernel and
a finite odd quotient. Apply this first to the preimage of
$\Lambda_o$ in $\mathcal P_2$ to obtain $\Delta_9$. Conjugation
embeds this complement into $\mathcal A_{\partial}^{\rm perm}$.
Its homology action is faithful, because the complete $Q_0$ action
is already faithful modulo two in D.34.

For a prescribed unit $\eta$, put $\lambda=N(\eta)$ and form the
integral homology map
$T=\lambda I_{\rm cusp}\oplus\eta_{\rm ell}$ using D.34's
canonical splitting. The central order-three deck map is the identity
on the cusp lattice and has the indicated CM action on the elliptic
lattice. The translation subgroup of order three permutes the cusp
classes and is the identity on the elliptic quotient. It follows that
$T$ centralizes the homology action of $\Delta_9$. Equation
(D.35.9) supplies an actual $a\in\mathcal A_{\partial}$ with
homology $T$.

Write $h$ for integral homology. In the preimage
$h^{-1}(h(\Delta_9))\subset\mathcal A_{\partial}^{\rm perm}$,
the two groups $\Delta_9$ and $a\Delta_9a^{-1}$ are complements
to the pro-two kernel of $h$. The preceding conjugacy proof gives
$l\in\ker h$ such that
$la\Delta_9(la)^{-1}=\Delta_9$. Because $T$ centralizes
$h(\Delta_9)$ and $h$ is injective on that finite subgroup,
$la$ centralizes it pointwise. The correction $l$ is pure with
multiplier one, so $la$ has precisely the required homology and
boundary powers. This proves the actual lifting assertion.

In (D.35.2), multiplying $1-2\zeta_3^2$ by its conjugate gives norm
seven, and expansion gives
$(1-2\zeta_3^2)^2/7=(-3-8\zeta_3^2)/7$.
Thus $N(\eta_*)=1$. Substitution of the matrix $Z$ gives
\[
 \eta_*(Z)=\begin{pmatrix}5/7&-8/7\\8/7&-3/7\end{pmatrix}
       \equiv\begin{pmatrix}3&8\\8&11\end{pmatrix}\pmod {16},
                                                               \tag{D.35.12}
\]
as asserted. This exact calculation only compares the compact Tate
component, so it does not identify the other D.33 coordinates.

It remains to verify every map in (D.35.3)--(D.35.5).
The subgroup $\Lambda_o$ is normal in $\Lambda$, and together with
$\bar s$ generates $\Lambda$. The pure peripheral subgroup is
normalized by $\bar s$, which permutes the three cusp labels. Hence
conjugation by $\bar s$ defines the displayed involution, and its
fixed subgroup is exactly the centralizer of all of $\Lambda$.
The odd-deck construction proves the asserted surjection from $C_o$
to the CM unit group. On the compact Tate lattice $s$ acts by $-I$,
so conjugation by it fixes the torus target. On the cusp factor all
maps in question are scalar, so their multipliers are also unchanged.

For the right arrow of (D.35.5), restrict to $K_2$ as in D.34 and
take the outer class. Its image lies in $C_o^{\iota}$. Conversely
choose a representative $a$ of any outer class there. The boundary
condition gives $a(c_0)=kc_0^{\lambda_2}k^{-1}$ for some $k\in K_2$.
Replacing $a$ by $\operatorname{Int}(k^{-1})a$ makes this equality
exact, without changing its outer class. D.34's two-way extension
map then gives an automorphism over $Q_0$ with
$t\mapsto t^{(\lambda_2,1)}$. Finally an element of the kernel is
$\operatorname{Int}(k)$ on $K_2$. Its homology forces
$\lambda_2=1$, and its exact boundary condition forces
$k\in C_{K_2}(c_0)=\langle c_0\rangle$. The free group's trivial
center identifies these inner automorphisms with $\mathbb Z_2$.
This proves exactness, including the actual full multiplier pair.

We now prove the stronger torus surjectivity (D.35.4a) arithmetically.
Write $k=\mathbb Q(\zeta_3)$, $\mathcal O=\mathbb Z_2[\zeta_3]$
and $U=1+2\mathcal O$. On $E_3:Y^2=X^3+1$, all three nonzero
two-torsion points are $k$-rational, with $X=-1,-\zeta_3,-\zeta_3^2$.
The CM automorphism $[\zeta_3]$ is also defined over $k$ and
permutes those three points cyclically. Its action modulo two has
irreducible polynomial $X^2+X+1$. Consequently $T_2E_3$ is free of
rank one over $\mathcal O$, and its $G_k$ image is a closed subgroup
$H\subset U$.

There is a small exact Frobenius computation proving $H=U$.
Direct counting gives
\[
 \begin{array}{c|c|c|c}
 p&\#E_3(\mathbb F_p)&a_p&\text{two CM Frobenius scalars}\\\hline
 7&12&-4&-1+2\zeta_3,\ -3-2\zeta_3\\
 13&12&2&3+4\zeta_3,\ -1-4\zeta_3.
 \end{array}                                         \tag{D.35.13}
\]
Both primes split in $k$. The characteristic polynomial of arithmetic
Frobenius on $T_2E_3$ is $X^2-a_pX+p$; its two roots in
$\mathcal O$ are the displayed scalars. This trace formula follows
from the degree of Frobenius and of $1-\mathrm{Frob}$ together with
the determinant-degree equality
\cite[II, Proposition 6.4; IV, Proposition 9.1 and the proof of
Theorem 9.4; V, Proposition 8.1]{dMilneEC2021}.
Conjugating an actual Frobenius by complex conjugation replaces its
CM scalar by its conjugate and interchanges the two primes over $p$.
Thus $H$ contains both listed roots for each row, independently of
the choice of the first prime above $p$.

Modulo eight, these four elements generate all sixteen elements of
$U/(1+8\mathcal O)$. Their multiplication is checked directly using
$(a+b\zeta_3)(c+d\zeta_3)
=(ac-bd)+(ad+bc-bd)\zeta_3$.
Moreover $1+8\mathcal O\subset U^2$: for $v\in8\mathcal O$,
the integral convergent series
$\exp(\tfrac12\log(1+v))\in1+4\mathcal O$ is a square root of
$1+v$. Hence $HU^2=U$. For completeness this implies $H=U$ without
an openness assumption: a proper closed subgroup of an abelian
pro-two group has a nontrivial finite two-group quotient, and that
quotient has a quotient of order two, which kills all squares. Such
a quotient contradicts $HU^2=U$.

The extension $k(\sqrt[3]{4})/k$ has degree three: a real cubic root
of four cannot lie in the quadratic field $k$. It is Galois since
$\zeta_3\in k$. Its intersection with the pro-two division field
$k(E_3[2^\infty])$ is $k$. Therefore the joint image of the actual
characters $(\rho_{E_3},\kappa_{4,3})$ is $U\times C_3$.
The exact twist (D.34.11) now gives
\[
          \rho_{\mathcal E_0}(G_k)
             =\mu_3\,U=\mathcal O^\times.            \tag{D.35.14}
\]
Every deck map in (D.34.2) is defined over $k$. Functoriality of the
geometric fundamental group consequently makes the actual outer
$G_k$ action commute with all of $\Lambda$, while all three rational
cusps are preserved with common multiplier $\chi_2$. This supplies
(D.35.4a). Its precise normalized lift is the one already constructed
in (D.35.5): inner conjugation by $K_2$ makes the $c_0$ image the
exact power, and the restriction/extension isomorphism of D.34 gives
the automorphism over $Q_0$ with full exponent $(\chi_2,1)$.
Inner conjugation by $K_2$ fixes the compact Tate action. Thus no
compact component changes in making this boundary normalization.

The involution itself also has a concrete quotient geometry. In the
model (D.34.1), $s$ is $z\mapsto1-z$ and fixes $w$.
Its quotient is the $w$ line, since
$z^2-z-w^3=0$. The three interior fixed points satisfy
$z=1/2$ and $1+4w^3=0$; the fourth fixed point is $P_\infty$.
The two removed points $P_0,P_1$ have common image $w=0$, and the
removed $P_\infty$ has image $w=\infty$. The quotient orbifold is
therefore the sphere with these two punctures and three order-two
cone points. Its fundamental group has the presentation
\[
 \langle q_1,q_2,q_3,h_0,h_\infty\mid
 q_1^2=q_2^2=q_3^2=1,\ q_1q_2q_3h_0h_\infty=1\rangle
       \simeq C_2*C_2*C_2*\mathbb Z.                 \tag{D.35.15}
\]
The covering character sends each $q_i$ and $h_\infty$ to the
nontrivial element of $C_2$, and $h_0$ to the identity. Its kernel is
the marked $K$, with the original cusp classes given by lifting these
puncture loops. Pro-two completion gives the exact preimage of
$\langle s\rangle$ in $\mathcal P_2$ already constructed in D.32:
$C_2\amalg_2C_2\amalg_2C_2\amalg_2\mathbb Z_2$.
This identifies the geometric involution and the previously used
Bass--Serre subgroup by their covering map, rather than just by their
abstract presentations.

Finally $D_9(K_2)$ is characteristic, so the group in (D.35.5)
acts on $R$. Its kernel $\langle\operatorname{Int}(c_0)\rangle$
maps onto the displayed procyclic inner subgroup of the finite target.
Quotienting gives exactly (D.35.5a). If a representative over $R$
differs from $\alpha_{33,R}$ by a power of this inner automorphism,
that power lifts to $\mathbb Z_2$ and is removed before reduction.
This proves the stated description of the actual fibre. The
single-torus arithmetic image (D.35.14) does not prescribe these
additional finite coordinates.

The accompanying program checks the two inverse basis changes, all
three boundary words, each displayed automorphism and inverse, the
four independent cusp shears, the generated group
$\operatorname{SL}_2(\mathbb Z/16)$ of order $3072$, the exact
norm-one lift (D.35.12), and every point count and unit-group
generation check in (D.35.13). The general two-adic and odd-deck arguments
above prove existence beyond these finite checks.
\end{proof}
